


%In this section, we review several existing methods that address aspects of openness or dynamic adaptability in multi-agent systems, including the techniques used in our investigations:  DGN~\cite{dgn}, DICG~\cite{dicg}, and GPL~\cite{gpl}, which focus on dynamic multi-agent coordination, and MOHITO~\cite{anilmohito}, which is specifically designed for task openness. We also discuss transformer-based models and briefly review multi-task, curriculum, and dynamic single-agent reinforcement learning approaches that adapt to evolving environments.

Graph Neural Networks (GNNs) \cite{gnn} model multi-agent interactions flexibly by representing agents as graph nodes and using message-passing to aggregate neighbor information, naturally handling variable group sizes but often limiting openness by imposing bounded assumptions. Several dynamic multi-agent approaches partially address related challenges. DGN~\cite{dgn} uses graph convolution with fixed-size neighbor communication for cooperation and can generalize to larger teams but restricts interaction through a predefined neighborhood and assumes static goals. DICG~\cite{dicg} learns dynamic coordination graphs for multi-agent coordination, handling large fixed teams well but requires known agent and task sets in advance, limiting true openness. GPL~\cite{gpl} tackles AO in ad hoc teamwork with graph neural nets and type inference for variable team sizes but trains only a single agent while others have fixed policies, thus not a fully adaptive multi-agent training approach. In contrast, MOHITO~\cite{anilmohito} is specifically designed for openness, formalizing task-open Markov games and employing hypergraph-based actor–critics to handle dynamic task–action sets, though it focuses solely on TO and assumes fixed agents and frames. Our investigations show that PLATO achieves superior performance compared to DGN, DICG, and MOHITO, with clear gains in harder openness setups, demonstrating the effectiveness of PLATO's pointing mechanism as an unbounded solution.
%transformers needs citations.

Also, while some may argue that transformers \cite{vaswani2017attention} naturally handle unbounded input sizes, imposing bounds is necessary to maintain training efficiency and scalability. That is, though transformers theoretically support variable-length inputs by attending over all tokens---making them attractive to OASYS, to enable efficient parallelization, \emph{inputs are typically padded to fixed lengths}.  This practice limits transformer-based applications in OASYS to a maximum number of agents or tasks. PLATO, on the other hand, circumvents this limitation via the attention mechanism and team-based statistics to handle variability in agents and tasks.  

Related fields such as multi-task reinforcement learning~\cite{zhang2021survey} and curriculum learning~\cite{narvekar2020curriculum} focus on adapting to varying tasks or progressively harder scenarios but typically assume a fixed number of entities and do not address AO and TO.  Also, they may suffer from catastrophic forgetting when learning new tasks~\cite{ribeiro2019multitask}, high prediction errors in volatile environments, and do not effectively handle openness due to their reliance on predefined task spaces or curriculum structures~\cite{kravchenko2024limitations}. %Dynamic single-agent reinforcement learning approaches can handle changing environment characteristics but remain limited to single-agent settings, exhibit poor sample efficiency in complex scenarios, and do not consider the multi-agent coordination challenges inherent in our open systems framework where agents, tasks, and interaction patterns can change simultaneously and unboundedly.




%\subsection{Related GNN-Based Approaches}\label{related:gnn}
%Graph neural networks (GNNs)~\cite{gnn} model agent interactions through message passing on graphs, naturally supporting variable team sizes. However, most GNN-based MARL methods still impose bounded adjacency or task sets, limiting full openness.

%GNNs \cite{gnn} offer a flexible framework to model multi-agent interactions in variable-sized groups by representing agents as nodes in a graph. In a GNN-based approach, edges (or an adjacency function) define which agents interact, and message-passing or convolution operations allow each agent to aggregate information from its neighbors. This graph-based representation naturally adapts when the team composition or interaction topology changes, without requiring a fixed input size. Several prior MARL approaches have attempted to handle aspects of openness, though each has limitations. In summary, they still impose bounds or assumptions that limit true openness. 

%For example, a standard multilayer perceptron (MLP) \cite{rumelhart1986learning_mlp} or Convolutional Neural Network (CNN) \cite{lecun2002gradient_cnn} requires a fixed-size input vector, so handling a variable number of entities often means padding the input to a maximum size and using masking to ignore the dummy entries or dynamically changing action spaces. This introduces an implicit upper bound on how many entities can be handled. Similarly, transformer-based \cite{wolf2020transformers} or attention-based \cite{vaswani2017attention} models input sequences are padded to a fixed length and masked. These workarounds mean the model is only trained up to a certain capacity – scenarios with more agents, task, or new frames beyond this will fall outside its training distribution. Prior multi-agent RL methods that appear to handle dynamic teams typically rely on such fixed-size placeholders. It has been observed that a changing number of agents yields a state representation of variable length, which causes traditional RL methods with fixed-length state vectors to perform poorly. In other words, without a specially designed architecture, the model either fails to process a different team size or must be tricked via padding/masking, often degrading performance in unseen or un-bounded conditions caused by openness.

%paragraph 4 DGN:
%DGN
%Graph Convolutional Reinforcement Learning (DGN) \cite{dgn} introduced a GNN-based method for multi-agent cooperation, where each agent shares parameters and communicates with a fixed number of neighbors (e.g., k-nearest agents) through graph convolution. This technique allows DGN to generalize to scenarios with more agents than seen in training (by exploiting relation kernels and neighbor-based message passing). However, DGN still imposes a structured communication graph (with a chosen neighborhood size k), effectively bounding the interactions and not truly allowing an unbounded number of agents beyond the predefined graph connectivity. Moreover, it assumes a static task goal and does not explicitly address agents entering or leaving during an episode.


%paragraph 5 MOHITO:

%MOHITO \cite{anilmohito} is one of the few approaches in MARL that explicitly tackles TO. It formalizes this setting with task-open Markov games (TaO-MG) and represents agent–task–action relations using interaction hypergraphs, which are converted into incidence graphs and processed by graph attention networks within an actor–critic framework. This enables reasoning over dynamically changing task–action sets without relying on a fixed task universe. MOHITO demonstrates effective performance in task-open domains such as ridesharing and wildfire, but it focuses solely on TO and does not address AO or FO.
%SHOULD I talk about the scalability of MOHITO and the perfomance (unbiasadli)? TODO

%paragraph 6 DICG:

%Deep Implicit Coordination Graphs (DICG) \cite{dicg} similarly leverage a graph-structured approach: DICG learns to infer a dynamic coordination graph among agents and uses a graph neural network to approximate the joint action-value function. This enables coordination in environments with many agents by learning which agents should interact for a given situation. However, DICG \cite{dicg} still presumes a fixed set of agents and tasks during training; it excels at coordinating large teams, but not at handling new agents or tasks appearing ad hoc. Any task openness in DICG’s experiments is limited to toggling the activity of known tasks (activating or deactivating predefined objectives), rather than introducing entirely new objectives. In other words, the set of possible tasks remains fixed, representing only a low level of task openness. DICG does not support unboundedness in the problem definition, since the universe of agents and tasks must be modeled in advance.

%paragraph 7 GPL:
%Graph-based Policy Learning (GPL) \cite{gpl} addresses agent openness in ad hoc teamwork \cite{stone2010ad}, where teammates with different fixed policies may enter or leave dynamically. GPL assumes that teammate types are fixed, implemented as heuristics or pretrained policies, and focuses on training a single learning agent rather than a full multi-agent system. Although it employs graph neural networks and coordination graphs to handle variable team sizes and is in principle unbounded in the number of teammates it can model, GPL is not a multi-agent reinforcement learning approach in the strict sense, since only one agent is being trained while others remain fixed. Its architecture combines type inference, an agent model to predict teammate actions, and a joint action value model to estimate their impact on the learner’s returns, to maximize the learner's single reward in the team.



%\subsection{Transformers}\label{sec:related}

%Transformers are a class of deep neural network architectures built on multi-head self-attention mechanisms, originally introduced for sequence modeling tasks like machine translation \cite{vaswani2017attention}. In a transformer, each element of the input can attend to every other element, enabling the model to capture rich dependencies without the sequential bottleneck of recurrent networks. This attention-based design allows transformers to process all input tokens in parallel and model long-range or pairwise relationships effectively. One major advantage over recurrent sequence-to-sequence models is this parallelization, which significantly accelerates training. In principle, transformers can accept variable-length sequences since their computations scale with the number of input tokens rather than requiring a fixed input size. This flexibility makes them attractive for MARL settings with dynamic agent or task sets, as the self-attention could, in theory, adapt to changing input sizes by attending over all present entities.

%However, in practice, efficient batching on GPUs/TPUs requires inputs to share a consistent tensor shape. To achieve this, implementations typically pad or truncate sequences to a uniform length and apply masks so that attention ignores dummy tokens \cite{wolf2020transformers}. While this enables fast parallelization, it introduces an implicit upper bound on the number of tokens the model can handle, wastes computation on padded positions, and constrains scalability. For instance, popular transformer networks set a maximum sequence length to control memory usage and quadratic runtime complexity, which becomes prohibitive as the number of entities grows \cite{tay2023efficient}.

%In the context of OASYS, this means transformer-based policies still assume a maximum number of agents or tasks, limiting their ability to handle true unbounded openness. Although adaptations such as set transformers avoid padding, they often sacrifice training efficiency \cite{lee2019set}. Overall, while theoretically flexible, transformers in practice remain constrained by padding, masking, and implicit upper bounds in most MARL applications with some form of openness.

 
